\section{Introduction}
Large language models can reproduce social stereotypes even when they perform well on conventional natural-language-processing benchmarks. In question answering, this creates two distinct risks: a model may rely on a stereotype when the context does not support a specific answer, or it may allow a stereotype to interfere with an answer that is explicitly supported by evidence. BBQ was introduced to separate these behaviors in a controlled multiple-choice setting \citep{parrish2022bbq}.

This study applies BBQ to a heterogeneous set of contemporary LLM configurations. The goal is not to build a general-purpose model leaderboard, but to examine three questions:
\begin{enumerate}[label=\textbf{RQ\arabic*:}]
    \item How reliably do the evaluated models distinguish ambiguity from evidence-supported cases?
    \item When models make non-\unknown{} predictions, how strongly do their answers align with the stereotype-target direction measured by BBQ?
    \item Does accuracy change when disambiguating evidence conflicts with the stereotype-aligned target?
\end{enumerate}

We evaluate the complete 58,492-example benchmark and report accuracy, directional bias, category-level behavior, and data-quality constraints together. This is important because a directional bias score near zero can coexist with poor uncertainty handling or poor evidence-following.
